\chapter{Glossary of Complexity Classes}
\label{app:glossary}

The classes used in the book, with one line definitions and the inclusions among them. Throughout, polynomial means polynomial in the input length $n$.

\begin{description}[style=nextline]
\item[$\Pc$] Languages decided by a deterministic algorithm in polynomial time.
\item[$\NP$] Languages with polynomial time verifiable certificates. $\Pc \subseteq \NP$.
\item[$\mathrm{co}\NP$] Complements of $\NP$ languages.
\item[$\ZPP$] Languages decided by a randomized algorithm that is always correct and runs in expected polynomial time. $\ZPP = \RP \cap \coRP$.
\item[$\RP$] Languages decided by a randomized polynomial time algorithm that never accepts a string outside the language and accepts each string inside with probability at least $\tfrac12$. $\Pc \subseteq \ZPP \subseteq \RP \subseteq \NP$.
\item[$\coRP$] Complements of $\RP$ languages. Primality was in $\coRP$ before 2002; polynomial identity testing is in $\coRP$.
\item[$\BPP$] Languages decided by a randomized polynomial time algorithm with two sided error at most $\tfrac13$. $\RP \subseteq \BPP$; $\BPP \subseteq \Sigma_2^p \cap \Pi_2^p$ (Sipser, G\'acs, Lautemann); $\BPP \subseteq \Ppoly$ (Adleman for $\RP$, Bennett and Gill for $\BPP$).
\item[$\mathrm{PP}$] Languages decided by a randomized polynomial time algorithm with acceptance probability greater than $\tfrac12$ on strings inside and at most $\tfrac12$ outside, with no gap. $\NP \subseteq \mathrm{PP}$ and $\Pc^{\mathrm{PP}} = \Pc^{\sharpP}$.
\item[$\sharpP$] Functions counting the certificates of an $\NP$ relation. The permanent, $\#\mathrm{SAT}$ and the queens completion count are $\sharpP$ complete.
\item[$\oplus\Pc$] Languages of the form $\{ x : f(x) \text{ odd} \}$ for $f \in \sharpP$. $\NP \subseteq \RP^{\oplus\Pc}$ (Valiant and Vazirani).
\item[$\PH$] The polynomial hierarchy, the union of $\Sigma_k^p$ over $k$, where $\Sigma_k^p$ is defined by $k$ alternating quantifiers with a polynomial time predicate. $\PH \subseteq \Pc^{\sharpP}$ (Toda).
\item[$\Ppoly$] Languages decided by polynomial size circuit families, not necessarily uniformly computable.
\item[$\mathrm{E}$, $\mathrm{EXP}$, $\mathrm{NEXP}$] Deterministic time $2^{O(n)}$, deterministic time $2^{\mathrm{poly}(n)}$, and nondeterministic time $2^{\mathrm{poly}(n)}$.
\item[FPRAS] A fully polynomial randomized approximation scheme: a randomized algorithm that, given $x$ and $\epsilon$, outputs a $(1 \pm \epsilon)$ approximation with probability at least $\tfrac34$ in time polynomial in $|x|$ and $1/\epsilon$. An FPRAS for a $\sharpP$ complete self reducible problem whose decision version is $\NP$ complete implies $\NP = \RP$.
\end{description}

The chain of inclusions that the book is about is
\[
\Pc \subseteq \ZPP \subseteq \RP \subseteq \NP \cap \BPP , \qquad \BPP \subseteq \Sigma_2^p \subseteq \PH \subseteq \Pc^{\sharpP} ,
\]
and the conjectures are $\Pc = \ZPP = \RP = \BPP$ and $\Pc \neq \NP$.

\chapter{Code}
\label{app:code}

The core Python routines behind the tables and figures of Chapters~\ref{ch:rare} and~\ref{ch:experiments}; the remaining scripts that produced the computed numbers of the book are in the \texttt{code/} directory, and the exact counts $Q(N)$ of Table~\ref{tab:counts} are taken from the literature (OEIS A000170; for $N = 27$, Preusser and Engelhardt~\cite{PreusserEngelhardt}), not produced by these scripts, which enumerate exhaustively only for small $N$. The listings are excerpts, not self contained programs: they assume the imports shown and the helper functions of \texttt{code/}. The energy function and its incremental update are the only problem specific parts; everything else is the vertical ladder.

\begin{footnotesize}
\begin{verbatim}
import math, random, itertools
import numpy as np
rng = random.Random(1)

def energy(p):
    n = len(p); c = 0
    for i in range(n):
        for j in range(i + 1, n):
            if abs(p[i] - p[j]) == j - i: c += 1
    return c

def delta_energy(p, i, j):
    n = len(p); old = 0; new = 0
    q = list(p); q[i], q[j] = q[j], q[i]
    for k in range(n):
        if k != i and k != j:
            if abs(p[i] - p[k]) == abs(i - k): old += 1
            if abs(p[j] - p[k]) == abs(j - k): old += 1
            if abs(q[i] - q[k]) == abs(i - k): new += 1
            if abs(q[j] - q[k]) == abs(j - k): new += 1
    if abs(p[i] - p[j]) == abs(i - j): old += 1
    if abs(q[i] - q[j]) == abs(i - j): new += 1
    return new - old

def levelset_chain(p, e, m, steps):
    n = len(p)
    for _ in range(steps):
        i, j = rng.sample(range(n), 2)
        d = delta_energy(p, i, j)
        if e + d <= m:
            p[i], p[j] = p[j], p[i]; e += d
    return p, e

def split_count(n, rho=0.3, M=400, sweeps=4):
    pop = []
    for _ in range(M):
        p = list(range(n)); rng.shuffle(p)
        pop.append((p, energy(p)))
    logZ = 0.0; m = 10**9; ladder = []
    while True:
        es = sorted(e for _, e in pop)
        mnew = es[int(rho * M) - 1]
        if mnew == m: mnew = m - 1
        if mnew < 0: mnew = 0
        frac = sum(1 for e in es if e <= mnew) / M
        logZ += math.log(frac); ladder.append((mnew, frac))
        if mnew == 0: break
        elite = [(list(p), e) for p, e in pop if e <= mnew]
        pop = []
        for i in range(M):
            p, e = elite[i % len(elite)]
            p, e = levelset_chain(list(p), e, mnew, sweeps * n)
            pop.append((p, e))
        m = mnew
    return math.exp(logZ) * math.factorial(n), ladder

def density_of_states(n):
    cnt = {}
    for p in itertools.permutations(range(n)):
        e = energy(p); cnt[e] = cnt.get(e, 0) + 1
    return cnt
\end{verbatim}
\end{footnotesize}

The bottom rung mixing experiment of Chapter~\ref{ch:experiments2} runs the same level set chain from a solution and records the first row column:

\begin{footnotesize}
\begin{verbatim}
def autocorrelation_time(x):
    x = np.array(x, float); x -= x.mean()
    v = x.var(); n = len(x)
    tau = 1.0
    for k in range(1, n // 4):
        c = np.dot(x[:-k], x[k:]) / ((n - k) * v)
        if c < 0.02: break
        tau += 2 * c
    return tau

def bottom_rung(n, steps):
    p = find_solution(n); e = 0; cols = []
    for _ in range(steps):
        i, j = rng.sample(range(n), 2); d = delta_energy(p, i, j)
        if e + d <= 1:
            p[i], p[j] = p[j], p[i]; e += d
        cols.append(p[0])
    return autocorrelation_time(cols)
\end{verbatim}
\end{footnotesize}

The connectivity computation of Chapter~\ref{ch:open} enumerates the level set $\{ S \le 1 \}$ by backtracking with at most one conflict, builds the swap graph on it, and finds components by breadth first search; its core is the backtracking enumerator:

\begin{footnotesize}
\begin{verbatim}
def enumerate_le1(n):
    out = []; x = [0] * n; used = [False] * n
    def rec(i, conf):
        if i == n:
            out.append(tuple(x)); return
        for c in range(n):
            if used[c]: continue
            k = sum(1 for j in range(i)
                    if abs(x[j] - c) == i - j)
            if conf + k > 1: continue
            x[i] = c; used[c] = True
            rec(i + 1, conf + k)
            used[c] = False
    rec(0, 0)
    return out
\end{verbatim}
\end{footnotesize}

The Wang Landau routine of Chapter~\ref{ch:experiments} uses the same proposal with the acceptance rule $\min\{1, g(S(x)) / g(S(x'))\}$ on the running estimate $g$, the update $\log g(S(x)) \leftarrow \log g(S(x)) + \log f$ after each step, and the halving of $\log f$ whenever the visit histogram is flat to within $20\%$.
